\section{Despliegue y confianza}

\subsection{Pila de evaluación y métricas de monitoreo}

Después del despliegue es necesario conectar el evidence de la etapa de entrenamiento con las métricas de producción:
\begin{itemize}
    \item Benchmark (HumanEval, LongBench) garantiza la capacidad;
    \item Guardrail (Hallucination, Forbidden topics) garantiza la seguridad;
    \item Ops metrics (Rejection rate, Latency P99) garantiza la experiencia.
\end{itemize}

\begin{table}[H]
\centering
\begin{tabular}{p{0.25\textwidth} p{0.55\textwidth}}
\toprule
Categoría & Métrica típica \\
\midrule
Benchmark & HumanEval, TruthfulQA, Domain-specific tests \\
Safety & Red team rejection, prompt-based filter precision \\
Ops & Rejection rate, Latency P99, Human override count \\
\bottomrule
\end{tabular}
\caption{Pila de monitoreo tras el despliegue del aprendizaje por imitación}
\end{table}

\subsection{Diapositivas y evidencia visual}

\begin{figure}[H]
\centering
\includegraphics[width=0.8\textwidth]{slide-02-05.jpg}
\caption{Slide 02-05 muestra el ciclo cerrado que conecta los datos de demostración con el monitoreo (data → policy → ops board).}
\end{figure}

\begin{figure}[H]
\centering
\includegraphics[width=0.8\textwidth]{slide-02-10.jpg}
\caption{Slide 02-10 combina el evaluation stack y el governance checkpoint en una sola matriz, y enfatiza el evidence-driven rollout.}
\end{figure}

\begin{importantbox}{El papel de los tableros de visualización}
Visualizar el evidence stacking como un mapa de calor en las diapositivas permite que PM, el equipo legal y SRE discutan los riesgos desde la misma perspectiva de datos, en lugar de depender de documentos dispersos.
\end{importantbox}

\subsection{Retroalimentación y gobernanza}

Se sigue el modelo Observe-Assess-Act-Document:
\begin{itemize}
    \item Observe: se registra el desempeño del rollout y el reward drift;
    \item Assess: se revisa el guardrail con la evidence matrix;
    \item Act: el human-in-the-loop activa el mitigation;
    \item Document: la decisión se registra en el runbook, para una auditoría posterior.
\end{itemize}

\begin{warningbox}{Documentar es más importante que corregir}
Un loop de retroalimentación sin runbook hace que al equipo le resulte difícil recordar por qué se hizo rollback en cierto prompt. Archit sugiere: cada mitigation debe ir acompañado de un drift snapshot, un human override log y un lesson learned.
\end{warningbox}

\subsection{Resumen de la sección}

La confianza en el despliegue proviene del ciclo cerrado entre data, policy y ops. Solo al convertir el evidence dashboard de las diapositivas en tablas reales de monitoreo y gobernanza se puede lograr que el sistema de aprendizaje por imitación se mantenga estable de manera continua en production.

